\begin{proof}
    Clearly $\sem{-}_H \circ \dagger$ and $\dagger \circ \sem{-}_H$ are both contravariant prop functors from $\CatHQC$ to a subcategory of $\CatHQC$.
    Then $\sem{-}_H \circ \dagger = \dagger \circ \sem{-}_H$ if and only if $\sem{-} \circ \dagger \circ i = \dagger \circ \sem{-}_H \circ i$.
    In other words, $\sem{U^\dagger}_H = \sem{U}_H^\dagger$ for all $U \in \Mor( \CatHQC )$ if and only if $\sem{g^\dagger}_H = \sem{g}_H^\dagger$ for all $g \in \Sigma_{\CatHQC}$.
    Let $g \in \Sigma_{\CatHQC}$.
    There are three cases to consider.
    \begin{enumerate}
    \item Assume $g \in \Sigma_{\textbf{Prim}}$.
          By definition, $g^\dagger = g$ and $\sem{g}_H$ is Hermitian and unitary.
          Since $\sem{g}_H$ is both Hermitian and unitary, then $\sem{g}_H = \sem{g}_H^\dagger$.
          In conclusion, $\sem{g^\dagger}_H = \sem{g}_H = \sem{g}_H^\dagger$.
    \item Assume $g \in \Sigma_{\textbf{Ctrl}}$ with $\dom( g ) = n$.
          Then by definition, $\sem{g^\dagger}_H = \sem{g \exparr (-1)}_H = \sem{g}_H^{-1}$.
          Then the following equation holds.
          \begin{align*}
              \sem{g}_H \sem{g^\dagger}_H
              &=
              \sem{g}_H^1 \sem{g}_H^{-1}
              \\
              &=
              \exp( \Log\sem{g}_H ) \cdot \exp( \Log\sem{g}_H (-1) )
              \\
              &=
              \exp( \Log\sem{g}_H + \Log\sem{g}_H (-1) )
              &
              ( \text{by~\cref{Eqn:SimulDiag}} )
              \\
              &=
              \exp( \mathbb{O}_{2^n} )
              =
              \mathbb{I}_{2^n}
          \end{align*}
          In other words, $\sem{g^\dagger}_H$ is the inverse to $\sem{g}_H$.
          Since $\sem{g}_H$ is unitary, then $\sem{g^\dagger}_H = \sem{g}_H^\dagger$.
    \item Assume $g \in \Sigma_{\textbf{Pow}}$.
          Then there exists a $U \in \Mor( \CatHQC )$ and an $\alpha \in \mathbb{R}$ such that $g = U \exparr \alpha$.
          Then $g^\dagger = U \exparr ( -\alpha )$.
          Then by definition, $\sem{g}_H = \sem{U}_H^\alpha$ and $\sem{g^\dagger} = \sem{U}_H^{-\alpha}$.
          Then the following equation holds.
          \begin{align*}
              \sem{g}_H \sem{g^\dagger}_H
              &=
              \sem{U}_H^\alpha \sem{U}_H^{-\alpha}
              \\
              &=
              \exp( \Log\sem{g}_H \alpha ) \cdot \exp( \Log\sem{g}_H (-\alpha) )
              \\
              &=
              \exp( \Log\sem{g}_H \alpha + \Log\sem{g}_H (-\alpha) )
              &
              ( \text{by~\cref{Eqn:SimulDiag}} )
              \\
              &=
              \exp( \mathbb{O}_{2^n} )
              =
              \mathbb{I}_{2^n}
          \end{align*}
          In other words, $\sem{g^\dagger}_H$ is the inverse to $\sem{g}_H$.
          Since $\sem{g}_H$ is unitary, then $\sem{g^\dagger}_H = \sem{g}_H^\dagger$.
    \end{enumerate}
    In each case, $\sem{g^\dagger}_H = \sem{g}_H^\dagger$.
    Since $g$ was arbitrary, then $\sem{g^\dagger}_H = \sem{g}_H^\dagger$ for all $g \in \Sigma_{\CatHQC}$.
    In conclusion, $\sem{U^\dagger}_H = \sem{U}_H^\dagger$ for all $U \in \Mor( \CatHQC )$.
\end{proof}
